% Oversells — dars ishlanmasi shabloni (XeLaTeX).
% Lesson plan template: lesson data, aims and a table of the lesson stages (compile with XeLaTeX).
\documentclass[11pt,a4paper]{article}

\usepackage[margin=2cm]{geometry}
\usepackage{fontspec}
\setmainfont{Noto Serif}
\setsansfont{Noto Sans}
\usepackage{babel}
\babelprovide[import=uz-Latn, main]{uzbek}
% Boʻgʻin koʻchirish: satr oxirida ham, boshida ham kamida 2 harf qoladi.
\def\uzbekhyphenmins{22}\def\uzbekcyrhyphenmins{22}
\setlength{\emergencystretch}{3em}
\usepackage{amsmath, amssymb}
\usepackage[table]{xcolor}
\usepackage{tabularx, array}
\usepackage{enumitem}
\setlist{nosep, leftmargin=1.4em}
\pagestyle{empty}

\definecolor{dars}{HTML}{1F6F5C}
\newcolumntype{L}{>{\raggedright\arraybackslash}X}
\renewcommand{\arraystretch}{1.35}

% Jadval sarlavhasi / a table heading cell
\newcommand\sarlavha[1]{\textcolor{white}{\bfseries #1}}
% Maʼlumot qatori / a row of the lesson data
\newcommand\malumot[2]{\textbf{#1} & #2 \\ \hline}

\begin{document}

\begin{center}
  {\sffamily\small\color{dars} TASDIQLAYMAN: \rule{3cm}{0.4pt} \quad «\rule{0.8cm}{0.4pt}» \rule{2.5cm}{0.4pt} 2025-yil}\\[1em]
  {\sffamily\LARGE\bfseries\color{dars} DARS ISHLANMASI}
\end{center}

\vspace{0.5em}
\noindent
\begin{tabularx}{\linewidth}{|>{\columncolor{dars!8}}p{4.2cm}|L|}
  \hline
  \malumot{Fan}{Algebra}
  \malumot{Sinf}{8-«A» sinf}
  \malumot{Sana}{\rule{3cm}{0.4pt}}
  \malumot{Mavzu}{Kvadrat tenglama ildizlari formulasi}
  \malumot{Dars turi}{Yangi bilim beruvchi dars}
  \malumot{Dars metodlari}{Muammoli savol, «Aqliy hujum», juftlikda ishlash, mustaqil ish}
  \malumot{Jihozlar}{Darslik, tarqatma varaqlar, proyektor, taqdimot}
  \malumot{Oʻqituvchi}{\rule{5cm}{0.4pt}}
\end{tabularx}

\subsection*{\color{dars}Dars maqsadlari}
\begin{description}[leftmargin=0pt]
  \item[Taʼlimiy:] oʻquvchilar kvadrat tenglama ildizlari formulasini
    keltirib chiqarishni va uni masalalar yechishda qoʻllashni oʻrganadi.
  \item[Tarbiyaviy:] aniqlik, mantiqiy fikrlash va jamoada ishlash
    koʻnikmalarini tarbiyalash.
  \item[Rivojlantiruvchi:] taqqoslash, umumlashtirish va oʻz fikrini asoslab
    aytish koʻnikmalarini rivojlantirish.
\end{description}

\subsection*{\color{dars}Kutilayotgan natijalar}
Oʻquvchi:
\begin{itemize}
  \item diskriminantni hisoblaydi va ildizlar sonini aniqlaydi:
    $D = b^2 - 4ac$;
  \item ildizlarni formula yordamida topadi:
    $x_{1,2} = \dfrac{-b \pm \sqrt{D}}{2a}$;
  \item yechimni Viyet teoremasi bilan tekshiradi.
\end{itemize}

\subsection*{\color{dars}Darsning borishi}
\noindent
\begin{tabularx}{\linewidth}{|c|>{\raggedright\arraybackslash}p{3.9cm}|L|c|}
  \hline
  \rowcolor{dars}\sarlavha{№} & \sarlavha{Bosqich} & \sarlavha{Oʻqituvchi va oʻquvchilar faoliyati} & \sarlavha{Vaqt} \\ \hline
  1 & Tashkiliy qism & Salomlashish, davomatni aniqlash, darsga tayyorgarlikni tekshirish. & 2 daq \\ \hline
  2 & Uy vazifasini tekshirish & 3--4 oʻquvchi doskada; qolganlar daftarni juftlikda tekshiradi. & 5 daq \\ \hline
  3 & Bilimlarni faollashtirish & Muammoli savol: $x^2 + 6x + 5 = 0$ tenglamani toʻla kvadrat ajratib yechish. & 5 daq \\ \hline
  4 & Yangi mavzu bayoni & Umumiy holda $ax^2 + bx + c = 0$ uchun formulani keltirib chiqarish; $D > 0$, $D = 0$, $D < 0$ hollari. & 15 daq \\ \hline
  5 & Mustahkamlash & Darslikdagi 112--115-mashqlar; juftlikda ishlash, natijalarni muhokama qilish. & 12 daq \\ \hline
  6 & Baholash & Faol oʻquvchilarni ragʻbatlantirish, baholarni izohlash. & 4 daq \\ \hline
  7 & Uy vazifasi & 116--119-mashqlar; qoʻshimcha: formulaning $b = 2k$ holidagi koʻrinishini topish. & 2 daq \\ \hline
  \multicolumn{3}{|r|}{\bfseries Jami} & \bfseries 45 daq \\ \hline
\end{tabularx}

\subsection*{\color{dars}Oʻqituvchining dars tahlili}
\vspace{0.3em}
\noindent\rule{\linewidth}{0.4pt}\\[1.2em]
\rule{\linewidth}{0.4pt}

\end{document}
